\section{枚举与剪枝}\label{ux679aux4e3eux4e0eux526aux679d}

\subsection{子集枚举}\label{ux5b50ux96c6ux679aux4e3e}

\passthrough{\lstinline!for (int s = mask; s; s = (s - 1) \& mask)!} 枚举 \passthrough{\lstinline!mask!} 的所有非空子集，总复杂度是 \passthrough{\lstinline!O(3\^n)!}（对所有 \passthrough{\lstinline!mask!} 求子集时）。

\begin{lstlisting}[language={C++}]
// 接口：enumerate_submasks(n)：按 mask 从小到大枚举每个 mask 的全部子掩码（含空集）；返回 {mask,sub} 序列。
// 参数：二进制集合使用的位数 n；要求 0<=n<31。
vector<pair<int, int>> enumerate_submasks(int n) {
    // 变量：result=按枚举顺序保存的 {mask,sub} 对。
    vector<pair<int, int>> result;
    // 变量：mask=当前子集/博弈状态的二进制掩码 mask。
    for (int mask = 0; mask < (1 << n); ++mask) {
        // 变量：sub=当前枚举的子掩码 sub。
        for (int sub = mask;; sub = (sub - 1) & mask) {
            result.push_back({mask, sub});
            if (sub == 0) break;
        }
    }
    return result;
}
\end{lstlisting}

\subsection{回溯剪枝}\label{ux56deux6eafux526aux679d}

固定顺序、提前判非法、记录当前最优上界、相同状态记忆化。对排列/组合先排序再跳过相同元素：\passthrough{\lstinline!if (i > start \&\& a[i] == a[i-1]) continue;!}。

\begin{lstlisting}[language={C++}]
// 接口：dfs(pos, sum)：枚举非负数组 a 中互不重复的组合并统计和为 target 的方案；累计和使用 __int128 避免加法溢出。
// 参数：当前 0-based 位置 pos；当前精确累计值 sum；要求 a 已排序且元素非负、answer 的方案计数可由 long long 表示。
void dfs(int pos, __int128 sum) {
    if (sum == target) { ++answer; return; }
    if (pos == n || sum > target) return;
    // 变量：i=当前循环下标 i。
    for (int i = pos; i < n; ++i) {
        if (i > pos && a[i] == a[i - 1]) continue;
        dfs(i + 1, sum + a[i]);
    }
}
\end{lstlisting}

\subsection{折半搜索（Meet-in-the-middle）}\label{ux6298ux534aux641cux7d22meet-in-the-middle}

\passthrough{\lstinline!n ≤ 40!}、每个元素选/不选且目标为和/异或/最值时，把数组分两半枚举，排序一侧后二分或双指针，复杂度 \passthrough{\lstinline!O(2\^(n/2) log 2\^(n/2))!}。

\begin{lstlisting}[language={C++}]
// 接口：gen(a, l, r, limit)：枚举 a 的半开区间 [l,r) 中不超过 limit 的子集和。
// 参数：非负数组 a；合法半开区间 [l,r)；非负上界 limit。超过上界的和在加法前剪枝。
vector<long long> gen(const vector<long long>& a, int l, int r,
                      long long limit) {
    // 变量：res=准备返回的结果容器/结果值 res。
    vector<long long> res{0};
    // 变量：i=当前循环下标 i。
    for (int i = l; i < r; ++i) {
        // 变量：old=修改前版本根/修改前容器大小 old。
        int old = (int)res.size();
        // 变量：j=内层循环下标 j。
        for (int j = 0; j < old; ++j)
            if (a[i] <= limit - res[j]) res.push_back(res[j] + a[i]);
    }
    return res;
}

// 接口：mitm_max_subset_sum(a, limit)：返回不超过 limit 的最大子集和；要求 a 中元素非负且 limit>=0。
// 参数：非负整数数组 a；子集和上界 limit。
long long mitm_max_subset_sum(const vector<long long>& a, long long limit) {
    // 变量：n=输入数组元素个数。
    int n = (int)a.size();
    // 变量：L=左半部分的全部子集和；R=右半部分的全部子集和。
    auto L = gen(a, 0, n / 2, limit), R = gen(a, n / 2, n, limit);
    sort(R.begin(), R.end());
    // 变量：best=目前找到的不超过 limit 的最大子集和。
    long long best = 0;
    // 变量：x=当前左半部分子集和。
    for (long long x : L) {
        if (x > limit) continue;
        // 变量：it=右半部分中第一个使 x+*it 超过 limit 的位置。
        auto it = upper_bound(R.begin(), R.end(), limit - x);
        if (it != R.begin()) best = max(best, x + *prev(it));
    }
    return best;
}
\end{lstlisting}

\section{前缀和、差分与扫描}\label{ux524dux7f00ux548cux5deeux5206ux4e0eux626bux63cf}

\subsection{一维前缀和}\label{ux4e00ux7ef4ux524dux7f00ux548c}

\passthrough{\lstinline!pre[i+1] = pre[i] + a[i]!}，区间 \passthrough{\lstinline![l,r]!} 和为 \passthrough{\lstinline!pre[r+1]-pre[l]!}。二维前缀和使用容斥：

\passthrough{\lstinline!S(x2,y2)-S(x1-1,y2)-S(x2,y1-1)+S(x1-1,y1-1)!}。

\subsection{一维差分}\label{ux4e00ux7ef4ux5deeux5206}

对 \passthrough{\lstinline![l,r]!} 加 \passthrough{\lstinline!v!}：\passthrough{\lstinline!d[l]+=v; d[r+1]-=v!}；最后累加 \passthrough{\lstinline!d!} 恢复数组。二维矩形加法在四个角标记，二维前缀恢复。

\begin{lstlisting}[language={C++}]
// 数值域：要求所有前缀和、区间修改累计量和修改后的 a[i] 均可由 long long 表示。
// 变量：pre=前缀和/前缀状态数组 pre；diff=当前差值/差分数组 diff。
vector<long long> pre(n + 1), diff(n + 2);
// 变量：i=当前循环下标 i。
for (int i = 1; i <= n; ++i) pre[i] = pre[i - 1] + a[i];
// 接口：range_sum(l, r)：供“一维差分”当前语句回调；返回值含义见调用处条件。
// 参数：闭区间左端点 l；闭区间右端点 r。
// 变量：range_sum=返回指定区间和的局部 lambda range_sum。
auto range_sum = [&](int l, int r) { return pre[r] - pre[l - 1]; };

// 接口：range_add(l, r, v)：供“一维差分”当前语句回调；返回值含义见调用处条件。
// 参数：闭区间左端点 l；闭区间右端点 r；顶点/状态编号 v。
auto range_add = [&](int l, int r, long long v) {
    diff[l] += v;
    diff[r + 1] -= v;
};
// 变量：l=当前事件/询问区间左端点 l；r=当前事件/询问区间右端点 r；v=当前相邻顶点/下一状态 v。
for (auto [l, r, v] : updates) range_add(l, r, v);
// 变量：cur=当前扫描位置的累计状态 cur。
long long cur = 0;
// 变量：i=当前循环下标 i。
for (int i = 1; i <= n; ++i) {
    cur += diff[i];
    a[i] += cur;
}
\end{lstlisting}

\subsection{扫描线思想}\label{ux626bux63cfux7ebfux601dux60f3}

把``开始/结束''事件排序，按坐标推进并维护当前集合。矩形面积、区间覆盖、会议室数量都可转成事件；需要动态维护覆盖长度时配合线段树。

\begin{lstlisting}[language={C++}]
vector<pair<int, int>> event; // {坐标, +1 开始 / -1 结束}
// 变量：l=当前事件/询问区间左端点 l；r=当前事件/询问区间右端点 r。
for (auto [l, r] : intervals)
    event.push_back({l, 1}), event.push_back({r, -1});
sort(event.begin(), event.end());
// 变量：active=扫描到当前位置时仍活跃的区间数 active；peak=扫描过程中出现过的最大活跃数量 peak。
int active = 0, peak = 0;
// 变量：x=当前输入值/查询值/坐标 x；delta=当前事件带来的计数增量 delta。
for (auto [x, delta] : event) {
    active += delta;
    peak = max(peak, active);
}
\end{lstlisting}

\section{排序、二分、双指针}\label{ux6392ux5e8fux4e8cux5206ux53ccux6307ux9488}

\subsection{二分查找}\label{ux4e8cux5206ux67e5ux627e}

\passthrough{\lstinline!lower\_bound!} 找第一个 \passthrough{\lstinline!>= x!}，\passthrough{\lstinline!upper\_bound!} 找第一个 \passthrough{\lstinline!> x!}。手写时推荐半开区间 \passthrough{\lstinline![l,r)!}：

\begin{lstlisting}[language={C++}]
// 变量：l=当前区间左端点 l；r=当前区间右端点 r。
int l = 0, r = n;
while (l < r) {
    // 变量：m=当前边数、操作数或第二维规模 m。
    int m = l + (r - l) / 2;
    if (a[m] < x) l = m + 1;
    else r = m;
}
// l 是第一个 a[l] >= x 的位置，可能等于 n
\end{lstlisting}

\subsection{二分答案}\label{ux4e8cux5206ux7b54ux6848}

当答案具有单调可行性 \passthrough{\lstinline!check(x)!} 时，在答案范围二分。先明确求``最小可行''还是``最大可行''，并把 \passthrough{\lstinline!check!} 写成无副作用函数。典型：最小最大分段和、最小速度、最大最小距离、满足容量的最少天数。

\begin{lstlisting}[language={C++}]
// 接口：minimize_max_segment_sum(a, k)：把非负数组分成至多 k 个非空连续段；答案可由 long long 表示时返回该值，否则返回 nullopt。
// 参数：非空非负数组 a；允许的最大分段数 k，要求 k>=1。
optional<long long> minimize_max_segment_sum(const vector<long long>& a, int k) {
    // 接口：check(limit)：判断能否把 a 分成至多 k 段且每段和不超过 limit。
    auto check = [&](__int128 limit) {
        // 变量：groups=当前贪心划出的连续段数。
        int groups = 1;
        // 变量：sum=当前连续段的元素和。
        __int128 sum = 0;
        // 变量：x=当前扫描的非负数组元素。
        for (long long x : a) {
            if (x > limit) return false;
            if (sum + x > limit) sum = x, ++groups;
            else sum += x;
        }
        return groups <= k;
    };
    // 变量：l=最大单元素值，亦为二分答案下界。
    __int128 l = *max_element(a.begin(), a.end());
    // 变量：r=全部元素之和，亦为二分答案上界。
    __int128 r = 0;
    // 变量：x=当前累加到宽整数总和中的非负数组元素。
    for (long long x : a) r += x;
    while (l < r) {
        // 变量：mid=当前二分判定的最大允许段和。
        __int128 mid = l + (r - l) / 2;
        if (check(mid)) r = mid;
        else l = mid + 1;
    }
    if (l > numeric_limits<long long>::max()) return nullopt;
    return (long long)l;
}
\end{lstlisting}

\subsection{双指针/滑动窗口}\label{ux53ccux6307ux9488ux6ed1ux52a8ux7a97ux53e3}

窗口右端只增不减，维护``当前窗口合法''不变量；若加入 \passthrough{\lstinline!r!} 使其非法，就移动 \passthrough{\lstinline!l!} 直到恢复。适用于正数和、至多/至少 \passthrough{\lstinline!k!} 个不同值、最长合法子串。含负数的和通常不能直接用双指针。

\begin{lstlisting}[language={C++}]
// 变量：cnt=计数数组/当前计数 cnt。
vector<int> cnt(alphabet);
// 变量：distinct=当前窗口内不同值的数量 distinct；left=滑动窗口左端点 left；answer=当前累计/最终答案 answer。
int distinct = 0, left = 0, answer = 0;
// 变量：right=滑动窗口右端点 right。
for (int right = 0; right < n; ++right) {
    if (cnt[s[right]]++ == 0) ++distinct;
    while (distinct > k) {
        if (--cnt[s[left++]] == 0) --distinct;
    }
    answer = max(answer, right - left + 1);
}
\end{lstlisting}

\subsection{三分与凸性}\label{ux4e09ux5206ux4e0eux51f8ux6027}

连续/离散单峰函数可三分；整数域更稳妥的做法是三分到很小区间后暴力。很多题可通过凸性改写为二分导数或斜率，不要盲目使用浮点三分。

\begin{lstlisting}[language={C++}]
// 数值域：要求 l<=r、r-l 可由 long long 表示且 r<LLONG_MAX；f(x) 的返回值也须可由 long long 表示。
while (r - l > 3) {
    // 变量：m1=第一个模数/第一部分规模 m1。
    long long m1 = l + (r - l) / 3;
    // 变量：m2=第二个模数/第二部分规模 m2。
    long long m2 = r - (r - l) / 3;
    if (f(m1) < f(m2)) l = m1 + 1; // 求最大值
    else r = m2 - 1;
}
// 变量：ans=当前累计答案 ans。
long long ans = LLONG_MIN;
// 变量：x=当前输入值/查询值/坐标 x。
for (long long x = l; x <= r; ++x) ans = max(ans, f(x));
\end{lstlisting}

\section{贪心}\label{ux8d2aux5fc3}

\subsection{典型证明模式}\label{ux5178ux578bux8bc1ux660eux6a21ux5f0f}

\begin{itemize}
\tightlist
\item
  区间调度：按右端点升序，选最早结束的可行区间；交换任意最优解的第一个区间即可证明。
\item
  最小生成树：割性质/环性质。
\item
  Huffman：每次合并最小的两个权值（最优前缀码、合并果子）。
\item
  任务排序：按截止时间、比值、边际收益排序，必须写出交换论证，不要凭直觉。
\end{itemize}

\begin{lstlisting}[language={C++}]
sort(intervals.begin(), intervals.end(),
     // 接口：匿名 lambda(x, y)：供“典型证明模式”当前语句回调；返回值含义见调用处条件。
     // 参数：待处理值或点/状态 x；待处理值或点/状态 y。
     [](auto &x, auto &y) { return x.second < y.second; });
// 变量：last=上一个已选择位置/值 last；chosen=“典型证明模式”这段代码中的临时量 chosen；值由本行初始化式确定。
int last = numeric_limits<int>::min(), chosen = 0;
// 变量：l=当前事件/询问区间左端点 l；r=当前事件/询问区间右端点 r。
for (auto [l, r] : intervals) {
    if (l >= last) last = r, ++chosen;
}
\end{lstlisting}

\subsection{反悔贪心}\label{ux53cdux6094ux8d2aux5fc3}

先按某个顺序加入方案，用堆保存可撤销元素；若当前超限，撤掉代价最大/收益最小者。常见于``最多完成多少任务''\,``带截止时间的作业选择''。

\begin{lstlisting}[language={C++}]
sort(jobs.begin(), jobs.end(),
     // 接口：匿名 lambda(x, y)：供“反悔贪心”当前语句回调；返回值含义见调用处条件。
     // 参数：待处理值或点/状态 x；待处理值或点/状态 y。
     [](auto &x, auto &y) { return x.deadline < y.deadline; });
// 变量：chosen=“反悔贪心”这段代码中的临时量 chosen；值由本行初始化式确定。
priority_queue<int> chosen;
// 变量：job=当前按贪心顺序处理的任务 job。
for (auto job : jobs) {
    chosen.push(job.cost);
    if ((int)chosen.size() > job.deadline) chosen.pop();
}
\end{lstlisting}

\subsection{差分约束式贪心}\label{ux5deeux5206ux7ea6ux675fux5f0fux8d2aux5fc3}

``每个位置至少/至多''\,``相邻差不超过''等限制有时可从左到右维护最小需求，再从右到左修正；若约束包含图上的环，转为最短路/最长路更可靠。

\begin{lstlisting}[language={C++}]
// x[v] <= x[u] + w；无解等价于存在可达负环
struct Edge { int v; long long w; }; // v=有向边终点，w=x[v]-x[u] 的 long long 上界。
// 变量：g=图/树的邻接表 g。
vector<vector<Edge>> g(n);
// 变量：d=当前距离、差值或覆盖增量 d。
vector<__int128> d(n, 0); // 宽整数势函数避免多条 long long 边连续松弛时溢出。
// 变量：inq=顶点当前是否在队列中的标记 inq；path_len=当前最短路所含边数，用于判负环。
vector<int> inq(n), path_len(n);
// 变量：q=当前 BFS/单调算法使用的队列 q。
queue<int> q;
// 变量：s=当前枚举的起点/源点 s。
for (int s = 0; s < n; ++s) q.push(s), inq[s] = 1;
// 变量：ok=当前条件是否仍成立 ok。
bool ok = true;
while (!q.empty() && ok) {
    // 变量：u=当前顶点/状态编号 u。
    int u = q.front(); q.pop(); inq[u] = 0;
    // 变量：v=当前边的终点 v；w=当前边权 w。
    for (auto [v, w] : g[u]) if (d[v] > d[u] + (__int128)w) {
        d[v] = d[u] + (__int128)w;
        path_len[v] = path_len[u] + 1;
        if (path_len[v] >= n) { ok = false; break; }
        if (!inq[v]) q.push(v), inq[v] = 1;
    }
}
\end{lstlisting}

\section{单调栈与单调队列}\label{ux5355ux8c03ux6808ux4e0eux5355ux8c03ux961fux5217}

\subsection{单调栈}\label{ux5355ux8c03ux6808}

维护栈内下标对应值单调。遇到更小/更大元素时弹栈，弹出元素的右侧第一个更小/更大位置被确定。应用：下一个更大元素、柱状图最大矩形、最大子矩形、贡献法统计子数组最小值。

\begin{lstlisting}[language={C++}]
// 变量：nextGreater=每个位置右侧第一个更大元素下标 nextGreater；st=当前集合或自动机/DP 状态 st。
vector<int> nextGreater(n, -1), st;
// 变量：i=当前循环下标 i。
for (int i = 0; i < n; ++i) {
    while (!st.empty() && a[st.back()] < a[i])
        nextGreater[st.back()] = i, st.pop_back();
    st.push_back(i);
}
\end{lstlisting}

\subsection{单调队列}\label{ux5355ux8c03ux961fux5217}

队首始终是窗口最优值，队尾删除不可能成为最优的元素。应用：滑动窗口最值、DP 转移 \passthrough{\lstinline!dp[i]=min(dp[j])+w(i)!}、有长度限制的最短/最长子数组。

\begin{lstlisting}[language={C++}]
// 变量：q=当前 BFS/单调算法使用的队列 q。
deque<int> q;
// 变量：i=当前循环下标 i。
for (int i = 0; i < n; ++i) {
    while (!q.empty() && q.front() <= i - k) q.pop_front();
    while (!q.empty() && a[q.back()] <= a[i]) q.pop_back();
    q.push_back(i);
    if (i >= k - 1) answer.push_back(a[q.front()]);
}
\end{lstlisting}

\section{离散化与坐标压缩}\label{ux79bbux6563ux5316ux4e0eux5750ux6807ux538bux7f29}

收集所有会出现的坐标，排序去重，用 \passthrough{\lstinline!lower\_bound!} 映射。若区间端点表示连续长度，通常还要把 \passthrough{\lstinline!x!} 与 \passthrough{\lstinline!x+1!} 或相邻端点间距保留，避免把长区间压成一个点。值域大、操作数小的线段树/树状数组必须先离散化。

\begin{lstlisting}[language={C++}]
// 变量：xs=离散化后排序去重的坐标 xs。
vector<long long> xs = coordinates;
sort(xs.begin(), xs.end());
xs.erase(unique(xs.begin(), xs.end()), xs.end());
// 接口：id(x)：供“离散化与坐标压缩”当前语句回调；返回值含义见调用处条件。
// 参数：待处理值或点/状态 x。
auto id = [&](long long x) {
    return int(lower_bound(xs.begin(), xs.end(), x) - xs.begin()) + 1;
};
\end{lstlisting}

\section{常见构造套路}\label{ux5e38ux89c1ux6784ux9020ux5957ux8def}

\begin{itemize}
\tightlist
\item
  先处理奇偶、余数、最大值/最小值等不变量。
\item
  需要构造排列时，优先考虑循环移位、两端放置、按奇偶分组。
\item
  证明``无解''时找模数、颜色、度数、连通性或总和不变量。
\item
  证明``总能构造''时给出明确操作顺序，并说明每一步都保持前置条件。
\end{itemize}

\begin{lstlisting}[language={C++}]
// 常见构造：把 1..n 按奇偶分组输出
// 变量：ans=当前累计答案 ans。
vector<int> ans;
// 变量：x=当前输入值/查询值/坐标 x。
for (int x = 1; x <= n; x += 2) ans.push_back(x);
// 变量：x=当前输入值/查询值/坐标 x。
for (int x = 2; x <= n; x += 2) ans.push_back(x);
\end{lstlisting}

\section{随机化与哈希技巧}\label{ux968fux673aux5316ux4e0eux54c8ux5e0cux6280ux5de7}

\begin{itemize}
\tightlist
\item
  随机 \passthrough{\lstinline!shuffle!} + 快速选择求期望线性第 k 小。
\item
  64 位随机哈希表示集合，支持多重集合差异的高概率比较；严谨题目用排序/Trie/计数。
\item
  Treap 随机优先级避免退化；不要把 \passthrough{\lstinline!rand()!} 作为高质量随机源，使用 \passthrough{\lstinline!mt19937\_64!}。
\end{itemize}

\begin{lstlisting}[language={C++}]
mt19937_64 rng(chrono::steady_clock::now().time_since_epoch().count());
shuffle(a.begin(), a.end(), rng);
nth_element(a.begin(), a.begin() + k, a.end());
// 变量：answer=当前累计/最终答案 answer。
int answer = a[k];
uint64_t random_hash = rng();
\end{lstlisting}

\section{位集优化}\label{ux4f4dux96c6ux4f18ux5316}

\passthrough{\lstinline!std::bitset!} 将 64/机器字并行，适合传递闭包（布尔矩阵）、集合交并、子集计数；\passthrough{\lstinline!bitset!} 大小需编译期常量，动态场景使用 \passthrough{\lstinline!vector<unsigned long long>!}。

\begin{lstlisting}[language={C++}]
// 变量：N=数组容量/预处理上界；实际合法下标以本节注释为准。
const int N = 2000;
// 接口：bitset_closure(reach, n)：原地计算 reach 前 n 行表示的传递闭包。
// 参数：reach 初始含自环和原图边；n 为实际顶点数且 n<=N。
void bitset_closure(array<bitset<N>, N>& reach, int n) {
    // 变量：k=当前允许作为中转点的顶点 k。
    for (int k = 0; k < n; ++k)
        // 变量：i=当前起点 i。
        for (int i = 0; i < n; ++i)
            if (reach[i][k]) reach[i] |= reach[k];
}
\end{lstlisting}

\section{快速小技巧}\label{ux5febux901fux5c0fux6280ux5de7}

\begin{itemize}
\tightlist
\item
  \passthrough{\lstinline!std::gcd!}、\passthrough{\lstinline!std::lcm!}（C++17）；后者仅在结果可由返回类型表示时使用。
\item
  \passthrough{\lstinline!\_\_builtin\_ctzll!} 只对非零数调用。
\item
  排序后用 \passthrough{\lstinline!unique!} 去重，差异数组用 \passthrough{\lstinline!adjacent\_difference!}。
\item
  大量模乘、叉积、斜率比较时先判断上界，必要时统一 \passthrough{\lstinline!\_\_int128!}。
\end{itemize}

\begin{lstlisting}[language={C++}]
// 接口：checked_lcm(a, b)：非负最小公倍数可由 long long 表示时返回该值，否则返回 nullopt；0 与任意数返回 0。
// 参数：任意 long long 整数 a、b，包括 LLONG_MIN 与 LLONG_MAX。
auto checked_lcm = [](long long a, long long b) -> optional<long long> {
    if (a == 0 || b == 0) return 0;
    // 先提升再取绝对值，避免对 LLONG_MIN 求负产生未定义行为。
    // 变量：x、y=提升到宽整数域后的两个操作数。
    __int128 x = a, y = b;
    if (x < 0) x = -x;
    if (y < 0) y = -y;
    // 变量：u、v=欧几里得算法的当前两项；g=绝对值最大公约数；value=精确最小公倍数。
    __int128 u = x, v = y;
    while (v != 0) {
        // 变量：r=本轮欧几里得算法的余数。
        __int128 r = u % v; u = v; v = r;
    }
    // 变量：g=绝对值最大公约数；value=宽整数域中的精确最小公倍数。
    __int128 g = u, value = x / g * y;
    if (value > numeric_limits<long long>::max()) return nullopt;
    return (long long)value;
};
// 变量：l=a 与 b 的非负最小公倍数；溢出时为 nullopt。
optional<long long> l = checked_lcm(a, b);
sort(v.begin(), v.end());
v.erase(unique(v.begin(), v.end()), v.end());
long long lowbit = x & -x; // x != 0
// 变量：safe_product=“快速小技巧”这段代码中的临时量 safe_product；值由本行初始化式确定。
long long safe_product = (long long)((__int128)x * y % MOD);
\end{lstlisting}
